\chapter[Preliminaries]{Mathematical Preliminaries}
\label{ch:preliminaries}

\paragraph{Learning outcomes.}
After this chapter the reader should be able to:
\begin{enumerate}
  \item distinguish a primitive metric from a composite vierbein
        bilinear, and write the conversion between \(g_{\mu\nu}\) and
        \(e^A_\mu\);
  \item vary the Einstein--Hilbert action to obtain
        \(G_{\mu\nu}=8\pi G\,T_{\mu\nu}\) and state the contracted
        Bianchi identity that forces
        \(\nabla^\mu T_{\mu\nu}=0\);
  \item expand a functional determinant to second Seeley--DeWitt order
        and recognize which coefficient multiplies \(R\);
  \item define a causal diamond, its modular Hamiltonian, and the
        first variation \(\delta S_D\);
  \item enter Chapter~3 already possessing every geometric sentence
        used in the matrix action and the later Ward identities.
\end{enumerate}

This chapter is a toolkit, not a treatise. Every section exists because
a later proof calls it by name. Results that are used as black boxes
are labeled as such. Results that NAMCG must derive, rather than
import, are flagged and postponed.

\section{Lorentzian manifolds and index conventions}
\label{sec:ch2-manifolds}

A spacetime in this book is a smooth four-manifold \(M\) equipped with
a non-degenerate symmetric tensor \(g_{\mu\nu}\) of signature
\((-,+,+,+)\). The inverse is \(g^{\mu\nu}\), and
\(g=\det(g_{\mu\nu})\). The volume form is
\(\sqrt{-g}\,\mathrm{d}^4x\).

Coordinate indices \(\mu,\nu,\rho,\sigma\) transform under
diffeomorphisms. Tangent-space (Lorentz) indices \(A,B,\ldots\)
transform under local \(\mathrm{SO}(3,1)\). Internal gauge indices
\(I,J,\ldots\) transform under the microscopic gauge group of the
condensate. These three alphabets are never mixed silently.

The Levi-Civita connection compatible with \(g_{\mu\nu}\) is
\begin{equation}
  \label{eq:christoffel}
  \Gamma^\rho_{\mu\nu}
  =\frac12 g^{\rho\sigma}
  \bigl(\partial_\mu g_{\nu\sigma}+\partial_\nu g_{\mu\sigma}
        -\partial_\sigma g_{\mu\nu}\bigr).
\end{equation}
The Riemann tensor, Ricci tensor, and Ricci scalar follow the
conventions
\begin{equation}
  \label{eq:riemann}
  \begin{aligned}
    R^\rho{}_{\sigma\mu\nu}
    &=\partial_\mu\Gamma^\rho_{\nu\sigma}
      -\partial_\nu\Gamma^\rho_{\mu\sigma}
      +\Gamma^\rho_{\mu\lambda}\Gamma^\lambda_{\nu\sigma}
      -\Gamma^\rho_{\nu\lambda}\Gamma^\lambda_{\mu\sigma},\\
    R_{\mu\nu}&=R^\rho{}_{\mu\rho\nu},\qquad
    R=g^{\mu\nu}R_{\mu\nu}.
  \end{aligned}
\end{equation}
The Einstein tensor is
\begin{equation}
  \label{eq:einstein-tensor}
  G_{\mu\nu}=R_{\mu\nu}-\frac12 g_{\mu\nu}R.
\end{equation}

\begin{remark}[Mostly-minus versus mostly-plus]
  Particle-physics texts often use \((+,-,-,-)\). Every energy
  condition and every sign in a heat-kernel coefficient is
  convention-dependent. If a formula in a cited paper looks off by a
  sign, check the signature before inventing new physics.
\end{remark}

\section{Vierbeine: primitive geometry versus composite geometry}
\label{sec:ch2-vierbein}

A vierbein (tetrad) is a local frame field \(e^A_\mu\) satisfying
\begin{equation}
  \label{eq:vierbein-def}
  g_{\mu\nu}=\eta_{AB}\,e^A_\mu e^B_\nu,\qquad
  \eta_{AB}=\mathrm{diag}(-1,+1,+1,+1).
\end{equation}
The inverse frame \(e^\mu_A\) obeys
\(e^A_\mu e^\mu_B=\delta^A_B\) and
\(e^A_\mu e^\nu_A=\delta^\mu_\nu\). Local Lorentz transformations act
as
\begin{equation}
  \label{eq:local-lorentz}
  e^A_\mu(x)\;\longmapsto\;
  \Lambda^A{}_B(x)\,e^B_\mu(x),
  \qquad
  \Lambda^T\eta\Lambda=\eta.
\end{equation}
The spin connection \(\omega_\mu{}^A{}_B\) is the unique metric-
compatible connection on the frame bundle (vanishing torsion in the
theories used here),
\begin{equation}
  \label{eq:spin-connection}
  \nabla_\mu e^A_\nu
  =\partial_\mu e^A_\nu-\Gamma^\rho_{\mu\nu}e^A_\rho
   +\omega_\mu{}^A{}_B e^B_\nu=0.
\end{equation}

\paragraph{Primitive metric.}
In textbook general relativity one takes \(g_{\mu\nu}\) as the
fundamental field. A vierbein is an optional rewriting, useful for
spinors.

\paragraph{Composite metric.}
In NAMCG the logical order is reversed. A microscopic matrix field
\(\Phi^I{}_A\) is fundamental. After condensation one defines
\begin{equation}
  \label{eq:composite-vierbein}
  e^A_\mu(x)\equiv\langle\partial_\mu\Phi^{IA}(x)\rangle
\end{equation}
and only then reconstructs \(g_{\mu\nu}\) by~\eqref{eq:vierbein-def}.
Equation~\eqref{eq:composite-vierbein} is not justified in this
chapter. It is the definition that Chapter~4 must earn.

The distinction matters for counting degrees of freedom. A primitive
symmetric tensor in four dimensions has ten components, of which
diffeomorphisms remove four and local Lorentz transformations are
absent because there is no independent frame. A primitive vierbein has
sixteen components; local Lorentz invariance removes six, diffeomorphisms
remove four, and six remain, matching the metric. When the vierbein is
composite, those sixteen numbers are not independent fields. They are
expectation values of derivatives of \(\Phi\). Gauge redundancies of
\(\Phi\) and of the emergent geometry must be matched, not double
counted. That matching is Chapter~10.

\section{Variation of the Einstein--Hilbert action}
\label{sec:ch2-EH}

The Einstein--Hilbert action with a cosmological term and a matter
action is
\begin{equation}
  \label{eq:EH-full}
  S[g]
  =\frac{1}{16\pi G}\int_M\mathrm{d}^4x\,\sqrt{-g}\,(R-2\Lambda)
   +S_{\mathrm{m}}[g,\Psi].
\end{equation}
Under a metric variation \(g_{\mu\nu}\to g_{\mu\nu}+\delta g_{\mu\nu}\)
one uses
\begin{equation}
  \label{eq:var-identities}
  \delta g^{\mu\nu}=-g^{\mu\alpha}g^{\nu\beta}\delta g_{\alpha\beta},\qquad
  \delta\sqrt{-g}=\frac12\sqrt{-g}\,g^{\mu\nu}\delta g_{\mu\nu},
\end{equation}
and the Palatini identity
\begin{equation}
  \label{eq:palatini}
  \delta R_{\mu\nu}
  =\nabla_\rho(\delta\Gamma^\rho_{\mu\nu})
   -\nabla_\nu(\delta\Gamma^\rho_{\mu\rho}).
\end{equation}
The covariant-derivative terms become a boundary integral. On a
manifold without boundary, or after adding the Gibbons--Hawking--York
term
\begin{equation}
  \label{eq:GHY}
  S_{\mathrm{GHY}}
  =\frac{1}{8\pi G}\int_{\partial M}\mathrm{d}^3x\,\sqrt{|h|}\,K,
\end{equation}
the bulk stationarity condition is
\begin{equation}
  \label{eq:EFE}
  G_{\mu\nu}+\Lambda g_{\mu\nu}=8\pi G\,T_{\mu\nu},
\end{equation}
where the matter stress tensor is defined by
\begin{equation}
  \label{eq:T-def}
  T_{\mu\nu}
  =-\frac{2}{\sqrt{-g}}
    \frac{\delta S_{\mathrm{m}}}{\delta g^{\mu\nu}}.
\end{equation}
The same definition will be used in Chapter~11 for the entropic
contribution \(\Tent_{\mu\nu}\), with \(S_{\mathrm{m}}\) replaced by
the entanglement functional.

\begin{remark}[Boundary terms are not optional in this book]
  Chapter~12 promotes a boundary entropy functional to a dynamical
  object. The combination \(K/8\pi G\) appearing in~\eqref{eq:GHY} is
  therefore not an afterthought. It is the gravitational piece of the
  HBEV boundary action.
\end{remark}

\section{Bianchi identities and conservation}
\label{sec:ch2-bianchi}

The Riemann tensor obeys
\begin{equation}
  \label{eq:bianchi}
  \nabla_{[\lambda}R_{\rho\sigma]\mu\nu}=0.
\end{equation}
Contracting twice yields the contracted Bianchi identity
\begin{equation}
  \label{eq:contracted-bianchi}
  \nabla^\mu G_{\mu\nu}=0.
\end{equation}
If the field equation is \(G_{\mu\nu}+\Lambda g_{\mu\nu}=8\pi G\,T_{\mu\nu}\),
then necessarily
\begin{equation}
  \label{eq:conservation}
  \nabla^\mu T_{\mu\nu}=0.
\end{equation}
This is not a courtesy to matter. It is the geometric content of
diffeomorphism invariance.

In an emergent theory the same identity becomes a constraint on the
microscopic construction. Suppose the total effective stress tensor is
split as
\begin{equation}
  \label{eq:T-split}
  T_{\mu\nu}^{\mathrm{total}}
  =T_{\mu\nu}^{\mathrm{matter}}
  +T_{\mu\nu}^{\mathrm{tensor}}
  +\Tent_{\mu\nu}.
\end{equation}
Geometry still requires
\(\nabla^\mu T_{\mu\nu}^{\mathrm{total}}=0\). It is not enough that
ordinary matter be conserved by itself if the extra pieces dump
momentum into the world. Chapter~11 exists to prove that the NAMCG
split satisfies~\eqref{eq:conservation}. Until that proof, any
inspiral ``drag'' is a phenomenological force and not a geometric
backreaction.

Noether's second theorem makes the same point in field language. If an
action is invariant under local diffeomorphisms generated by an
arbitrary vector field \(\xi^\mu\), the Euler--Lagrange tensor
\(\mathcal{E}^{\mu\nu}\) of the metric obeys
\(\nabla_\mu\mathcal{E}^{\mu\nu}=0\) identically, off shell as an
identity of the tensor densities involved, and on shell as conservation
of \(T^{\mu\nu}\). Mass terms that break
\(\delta_\xi h_{\mu\nu}=\nabla_\mu\xi_\nu+\nabla_\nu\xi_\mu\)
destroy that identity. That observation is the skeleton of Chapter~7.

\section{Gauge covariant derivatives and a scalar warmup}
\label{sec:ch2-gauge}

Let \(A_\mu=A_\mu^a T^a\) be a connection on a principal bundle with
compact gauge group, generators \(T^a\) anti-Hermitian in the
conventions of this book. The covariant derivative on a field \(\Phi\)
in a representation \(\rho\) is
\begin{equation}
  \label{eq:covdev}
  D_\mu\Phi=\partial_\mu\Phi-iA_\mu\Phi
\end{equation}
(up to the representation matrices, restored whenever the index
structure requires them). The field strength is
\begin{equation}
  \label{eq:fieldstrength}
  F_{\mu\nu}=\partial_\mu A_\nu-\partial_\nu A_\mu-i[A_\mu,A_\nu].
\end{equation}

A complex scalar with Mexican-hat potential,
\begin{equation}
  \label{eq:abelian-higgs}
  S[\phi]
  =\int\mathrm{d}^4x
   \Bigl(
     |D_\mu\phi|^2
     -\frac{\lambda}{4}\bigl(|\phi|^2-v^2\bigr)^2
   \Bigr),
\end{equation}
condenses at \(|\langle\phi\rangle|=v\). The phase is eaten by \(A_\mu\)
(Abelian Higgs / Stueckelberg mechanism), and the physical spectrum
contains a massive vector and a massive radial scalar. There is no
graviton.

This warmup is here for two reasons. First, the NAMCG potential in
Chapter~3 is the non-Abelian matrix analogue of~\eqref{eq:abelian-higgs}.
Second, the failure mode is already visible: a single complex scalar
has one Goldstone phase. After it is eaten, nothing remains that can
be identified with a symmetric trace-free tensor \(T_{\mu\nu}\).

\section{One-loop effective actions and the heat kernel}
\label{sec:ch2-heatkernel}

Let \(\mathcal{O}\) be a second-order elliptic (or hyperbolic, after
Wick rotation) differential operator on a bundle over \(M\), typically
of Laplace type,
\begin{equation}
  \label{eq:laplace-type}
  \mathcal{O}=-g^{\mu\nu}D_\mu D_\nu+E,
\end{equation}
with bundle endomorphism \(E\). The one-loop effective action obtained
by integrating out a fluctuation \(\eta\) is
\begin{equation}
  \label{eq:one-loop}
  W[g]
  =\frac12\mathrm{Tr}\,\ln\mathcal{O}
  =-\frac12\int_0^\infty\frac{\mathrm{d}s}{s}\,\mathrm{Tr}\,e^{-s\mathcal{O}},
\end{equation}
up to local counterterms and the usual \(i\varepsilon\) / Wick-rotation
caveats. The operator \(e^{-s\mathcal{O}}\) is the heat kernel. Its
diagonal admits the Seeley--DeWitt expansion
\begin{equation}
  \label{eq:SDW}
  \langle x|e^{-s\mathcal{O}}|x\rangle
  \sim
  \frac{1}{(4\pi s)^{2}}
  \bigl(
    a_0(x)+a_1(x)s+a_2(x)s^2+\cdots
  \bigr)
  \qquad(s\to 0^+),
\end{equation}
in four dimensions. For a Laplace-type operator the first coefficients
are
\begin{equation}
  \label{eq:SDW-coeffs}
  \begin{aligned}
    a_0&=1,\\
    a_1&=\frac16 R-E,\\
    a_2&=\frac{1}{180}\bigl(
           R_{\mu\nu\rho\sigma}R^{\mu\nu\rho\sigma}
          -R_{\mu\nu}R^{\mu\nu}
          +\tfrac72 R^2
         \bigr)
         +\frac12 E^2
         -\frac16 RE
         +\cdots,
  \end{aligned}
\end{equation}
where the omitted terms involve derivatives of \(E\) and the bundle
curvature. The precise tensor combination in \(a_2\) will be written
in full in Appendix~A when it is needed for ghost analysis. What
Chapter~8 requires is only this fact:

\begin{proposition}[Which coefficient produces Newton’s constant]
  \label{prop:a1-G}
  In the derivative expansion of \(W[g]\), the term linear in
  curvature and quadratic in the ultraviolet cutoff comes from
  \(a_1\). After a cutoff \(\LUV\) is imposed on the proper-time
  integral~\eqref{eq:one-loop}, one obtains a contribution of the
  schematic form
  \begin{equation}
    \label{eq:W-R}
    W[g]
    \supset
    \int\mathrm{d}^4x\,\sqrt{-g}
    \bigl(
      c_0\LUV^4+c_2\LUV^2 R+\cdots
    \bigr).
  \end{equation}
  The coefficient \(c_2\) is the microscopic ancestor of
  \(1/(16\pi\Geff)\). The coefficient \(c_0\) is the microscopic
  ancestor of the cosmological-constant problem.
\end{proposition}

\begin{proof}[Proof sketch]
  Insert~\eqref{eq:SDW} into~\eqref{eq:one-loop} and cut the
  proper-time integral at \(s\sim\LUV^{-2}\). The \(a_0/s^3\) piece
  produces \(\LUV^4\). The \(a_1/s^2\) piece produces \(\LUV^2\) times
  \(\tfrac16 R-E\). After the endomorphism \(E\) is evaluated on the
  condensate background, the curvature term remains and is absorbed
  into the Einstein--Hilbert prefactor. Details and the logarithmic
  \(v^2\ln(\LUV/\mu)\) piece are Chapter~8.
\end{proof}

Two warnings, both of which the later text will spend chapters
respecting.

\begin{enumerate}
  \item The quartic term \(c_0\LUV^4\) is not small. Declaring
        \(\Lambda_{\mathrm{eff}}\approx 10^{-120}\rho_{\mathrm{Pl}}\)
        without a mechanism is not a calculation. Chapter~12 is that
        mechanism, or it is a failure.
  \item If the fluctuation operator is fourth order, as in the
        quadratic action of Chapter~9, the heat-kernel calculus
        changes and Ostrogradsky poles appear. One cannot quote
        \(a_1\) from a second-order Laplacian and then use a
        \(\Box^2\) kinetic term without further work.
\end{enumerate}

\section{Causal diamonds}
\label{sec:ch2-diamonds}

Let \(p\) and \(q\) be events in \(M\) with \(q\) in the chronological
future of \(p\). The causal diamond (Alexandrov set)
\begin{equation}
  \label{eq:diamond}
  D(p,q)
  =J^+(p)\cap J^-(q)
\end{equation}
is the intersection of the causal future of \(p\) with the causal past
of \(q\). Its boundary \(\partial D\) consists of two null cones
joined at a spacelike 2-surface \(\mathcal{B}\) (the edge, or
bifurcation surface in the small-diamond limit).

A \emph{small} causal diamond is one whose proper size is short
compared with the local curvature radius, so that the metric inside
\(D\) may be written in Riemann normal coordinates
\begin{equation}
  \label{eq:RNC}
  g_{\mu\nu}(x)
  =\eta_{\mu\nu}
   -\frac13 R_{\mu\rho\nu\sigma}(0)x^\rho x^\sigma
   +O(x^3).
\end{equation}
Jacobson's thermodynamic derivation of Einstein's equation uses
precisely such diamonds: the matter flux across \(\partial D\) is
compared with the variation of entanglement entropy of the diamond
interior.

\emph{Dyadic} diamonds \(D_{j,n}\) are a nested family obtained by
repeated bisection of a reference diamond. They are bookkeeping. The
physics is in the statement that first variations of entropy vanish
for arbitrary small deformations of \(\partial D\), not in the dyadic
label.

\begin{definition}[Diamond entanglement functional]
  \label{def:SD}
  Let \(\rho_D=\mathrm{Tr}_{M\setminus D}\rho\) be the reduced state
  of the microscopic theory on the diamond interior. The von Neumann
  entropy
  \begin{equation}
    \label{eq:vonneumann}
    S_D=-\mathrm{Tr}(\rho_D\ln\rho_D)
  \end{equation}
  is the diamond entanglement functional. The modular Hamiltonian is
  \begin{equation}
    \label{eq:modular}
    H_D=-\ln\rho_D,
  \end{equation}
  so that \(S_D=\langle H_D\rangle\) up to the additive constant
  \(\mathrm{Tr}(\rho_D\ln\rho_D)+\langle\ln\rho_D\rangle=0\) written
  in that form.
\end{definition}

For a free field in Rindler space, \(H_D\) is proportional to the
boost generator (Bisognano--Wichmann). For a small diamond in a
curved background the same statement holds approximately, with
Unruh temperature
\begin{equation}
  \label{eq:unruh}
  T_{\mathrm{Unruh}}=\frac{\kappa}{2\pi},
\end{equation}
where \(\kappa\) is the surface gravity of \(\partial D\) as seen by
the conformal Killing vector that fixes \(\mathcal{B}\).

This book does \emph{not} take Jacobson's equation of state as an
axiom imported from thermodynamics of horizons. Chapter~6 derives
\(\delta S_D=0\) from the microscopic matrix theory, or states
clearly where an assumption remains. The definitions above are
only the language in which that derivation is written.

\section{Mutual information}
\label{sec:ch2-MI}

For two regions \(A\) and \(B\),
\begin{equation}
  \label{eq:MI}
  I(A:B)=S(A)+S(B)-S(A\cup B).
\end{equation}
Mutual information is non-negative, vanishes if and only if
\(\rho_{AB}=\rho_A\otimes\rho_B\) (for finite-dimensional algebras),
and is invariant under local unitaries on \(A\) or \(B\). In NAMCG the
macroscopic potential
\begin{equation}
  \label{eq:Uent}
  U_{\mathrm{ent}}
  =-\lambda\,I(A:B)
  =-\lambda\bigl[S(A)+S(B)-S(A\cup B)\bigr]
\end{equation}
is the interaction energy assigned to two condensed cores \(A\) and
\(B\). The sign is chosen so that positive mutual information lowers
the potential: entanglement is attractive in the effective
description.

The variation of \(U_{\mathrm{ent}}\) with respect to the worldline of
a core produces the drag acceleration of Chapter~11. Because
\(I(A:B)\) is built from diffeomorphism scalars \(S(\,\cdot\,)\), that
variation is not free to point in an arbitrary non-conservative
direction. Making this remark precise is the content of
Theorem~2 in the ultraviolet note, rewritten as a textbook proof in
Chapter~11.

\section{What this chapter deliberately omits}
\label{sec:ch2-omits}

The following topics are real, and they are not treated here.

\begin{itemize}
  \item A complete development of spinor bundles and the Dirac
        operator. NAMCG as formulated in the source notes is bosonic
        at the microscopic level. Fermionic matter can be coupled to
        the emergent vierbein later; it is not needed to obtain
        \(G_{\mu\nu}\).
  \item Full renormalization of quantum gravity as a fundamental
        metric theory. The metric is not fundamental here. The
        ultraviolet object is the matrix condensate, with cutoff
        \(\LUV\).
  \item A self-contained treatise on Tomita--Takesaki modular theory.
        Definition~\ref{def:SD} and the replica formula
        \(S=(1-n\partial_n)\ln\mathrm{Tr}\,\rho^n|_{n=1}\)
        are sufficient for Chapter~6. Appendix~B expands them.
  \item Numerical relativity. The integrators of Chapter~15 are
        ordinary-differential toy models, announced as such.
\end{itemize}

If a later proof appears to need a tool that is not in this chapter,
the defect is in the later proof or in this chapter, and should be
repaired rather than papered over with a citation.

\section{Exercises}
\label{sec:ch2-exercises}

\begin{enumerate}
  \item \textbf{Volume variation.}
    Derive both identities in~\eqref{eq:var-identities} from
    \(\delta g=g\,g^{\mu\nu}\delta g_{\mu\nu}\) and
    \(g^{\mu\alpha}g_{\alpha\nu}=\delta^\mu_\nu\).
  \item \textbf{Einstein tensor from variation.}
    Starting from~\eqref{eq:EH-full} with \(\Lambda=0\) and no
    boundary, obtain \(G_{\mu\nu}=8\pi G\,T_{\mu\nu}\). Keep the
    Palatini terms long enough to see them become a total derivative.
  \item \textbf{Contracted Bianchi.}
    Contract~\eqref{eq:bianchi} twice and obtain
    \(\nabla^\mu G_{\mu\nu}=0\). Do not use the field equations.
  \item \textbf{Frame versus metric degrees of freedom.}
    Count local components of \(e^A_\mu\) and of \(g_{\mu\nu}\).
    Subtract local Lorentz transformations and diffeomorphisms. Confirm
    that both descriptions have six physical components off shell in
    the usual sense of geometric gravity, before any equations of
    motion.
  \item \textbf{Composite-frame warning.}
    Suppose \(e^A_\mu=\langle\partial_\mu\Phi^{IA}\rangle\) and
    \(\Phi\) is subject to a global internal rotation
    \(\Phi\mapsto U\Phi\). What happens to \(e^A_\mu\)? What additional
    structure is required before this rotation can be identified with
    a local Lorentz transformation of the emergent frame? (The answer
    is Chapter~10; write the question cleanly now.)
  \item \textbf{Heat-kernel dimensions.}
    In four dimensions, assign canonical mass dimensions to
    \(a_0,a_1,a_2\) and to the proper-time variable \(s\). Show that
    \(a_0\LUV^4\) and \(a_1\LUV^2\) are the only two local terms
    generated by~\eqref{eq:SDW} that can appear as relevant or
    marginal couplings of mass dimension four or less, aside from
    logarithms.
  \item \textbf{Small diamond.}
    In Riemann normal coordinates~\eqref{eq:RNC}, compute the
    spatial 3-volume of a small diamond of radius \(\ell\) to order
    \(\ell^5 R\). This is the geometric ingredient in Jacobson's
    original argument.
  \item \textbf{Mutual information positivity.}
    Prove \(I(A:B)\ge 0\) from strong subadditivity, or from relative
    entropy \(S(\rho_{AB}\|\rho_A\otimes\rho_B)\). Why does the sign
    in~\eqref{eq:Uent} then correspond to an attractive effective
    potential?
\end{enumerate}

\section{Bridge to Chapter~3}
\label{sec:ch2-bridge}

The language is now in place. A metric may be primitive or composite.
If it is composite, it is built from a vierbein, and the vierbein must
come from a field that carries a tangent index. The Einstein tensor is
the Euler--Lagrange tensor of~\eqref{eq:EH-full}, and Bianchi forces
conservation. Newton's constant, if it is emergent, lives in the
\(a_1\) heat-kernel coefficient. Entanglement lives on causal
diamonds.

What is missing is the field. Chapter~3 introduces the non-Abelian
matrix order parameter \(\Phi^I{}_A\), the matrix-coordinate algebra
\([\hat X^\mu,\hat X^\nu]=i\Theta^{\mu\nu}\), and the microscopic
action whose condensation will produce the vierbein
of~\eqref{eq:composite-vierbein}. After that chapter, geometry is no
longer an independent ingredient. It is an output.